\subsubsection{Polar curvature in higher dimension}
\label{subsec:polar-first-current-obstruction}

The first curvature current no longer records all logarithmic poles in higher dimension. Its restriction to the pole set can vanish even when the model has nontrivial multiplier ideals. Thus the curve metric cannot be extended by replacing the polar measure with only this current restriction.

Let $X$ be a connected compact complex manifold of dimension $n\geq2$, and let $\omega$ be Hermitian. For $p\in\PSH(X,\omega)$ write
\[
 \operatorname{Pol}_1(p)
      =\mathbf1_{\{p=-\infty\}}(\omega+\ddc p),
\]
where restriction is taken on the coefficient measures of the positive current.

\Needspace{17\baselineskip}
\begin{proposition}\label{prop:isolated-pole-polar-current-blindness}\index{Model potential!polar current obstruction}\index{Multiplier ideal!isolated pole}
For every $x\in X$, there is $a_0>0$ such that each rational $a\in(0,a_0)$ is realized by an analytic ideal model $p_a$ with pole set $\{x\}$ and
\[
 p_a=a\log|z|^2+O(1)\quad\text{near }x,\qquad
 \operatorname{Pol}_1(p_a)=0.
\]
The models $p_a$ are pairwise ordinarily inequivalent and have different all-weight multiplier-ideal data, whereas $\operatorname{Pol}_1(0)=0$ also. More precisely, with $\mathfrak m_x$ the maximal ideal at $x$,
\begin{equation}\label{eq:isolated-pole-model-ideals}
 \mathcal I(tp_a)_x
       =\mathfrak m_x^{\max\{\lfloor ta\rfloor-n+1,0\}}
       \qquad(t>0).
\end{equation}
\end{proposition}
\begin{proof}
Choose the cutoff potentials $u_a$ of \cref{lem:small-isolated-log-poles} for the single point $x$, taking $a_0$ smaller than the reciprocal of its curvature-loss constant. Put $p_a=P_\omega[u_a]$. By \cref{prop:analytic-envelope-type}, this is an analytic ideal model ordinarily equivalent to $u_a$. It is locally bounded away from $x$, and its only pole has the displayed logarithmic coefficient. Distinct coefficients give unbounded differences near $x$. Formula \eqref{eq:isolated-pole-model-ideals} follows from \cref{lem:smooth-center-multiplier-ideal} with codimension $n$; its exponents distinguish two coefficients for sufficiently large weights.

It remains to check that the current has no coefficient mass at $x$. In coordinates centered at $x$, let $\beta=\ddc|z|^2$ and let
\[
 \nu=(\omega+\ddc p_a)\wedge\beta^{n-1}.
\]
This is a positive measure. Choose smooth cutoffs $0\leq\chi_\varepsilon\leq1$, equal to one on the ball of radius $\varepsilon$ and supported in the ball of radius $2\varepsilon$, with $-C\varepsilon^{-2}\beta\leq\ddc\chi_\varepsilon\leq C\varepsilon^{-2}\beta$. Their Hessians are supported on the intervening annulus. Since $\beta$ is closed, distributional integration by parts gives
\[
 \int\chi_\varepsilon\,\mathrm d\nu
 =\int\chi_\varepsilon\omega\wedge\beta^{n-1}
       +\int p_a\,\ddc\chi_\varepsilon\wedge\beta^{n-1}.
\]
The logarithmic bound on the annulus and its Euclidean volume therefore give
\[
 0\leq\nu(\{x\})\leq\int\chi_\varepsilon\,\mathrm d\nu
 \leq C_a\bigl(\varepsilon^{2n}
        +\varepsilon^{2n-2}(1+|\log\varepsilon|)\bigr)
       \longrightarrow0.
\]
Here $n\geq2$ is essential. Positivity bounds the coefficient measures of $\omega+\ddc p_a$ by its trace measure $\nu$. They consequently have no mass at $x$. Restriction to its pole set is zero, as asserted. The zero potential is an ordinary and ideal model because $\omega>0$, and its pole set is empty.
\end{proof}

Any comparison determined only by the fixed background and $\operatorname{Pol}_1(p)$ therefore fails to separate these analytic ideal models. Higher-codimensional singularity data must enter a complete higher-dimensional topology.
